\begin{theorem}[Regular representation multiplicities]
    \label{thm:peps_regular_fourier_multiplicities}
    \lean{TNLean.PEPS.leftRegularEuclideanMatrixEquiv,
        TNLean.PEPS.exists_unitary_leftRegular_matrix_blocks}
    \leanok
    \uses{thm:peps_unitary_character_blocks, thm:peps_semiregular_equivalence}
    The actual left-regular permutation representation of a finite group
    admits the orthonormal decomposition
    (\ref{eq:peps_unitary_irreducible_rows}), with $m_i=d_i$ for every
    irreducible sector. This is the regular decomposition used in
    \cite[Section~7]{Schuch2010PEPS}.
\end{theorem}
\begin{proof}\leanok
    The group basis identifies the group-algebra regular representation
    with the permutation matrices on the standard Hilbert space.
    Their characters agree. The multiplicity of a character $\chi$ in
    the regular representation is $\chi(1)$. For an irreducible row,
    $\chi(1)=d_i$, whereas the orthonormal row decomposition gives
    multiplicity $m_i$. Thus $m_i=d_i$.
\end{proof}

\begin{theorem}[Unitary Fourier decomposition of the regular matrices]
    \label{thm:peps_regular_unitary_fourier}
    \lean{TNLean.PEPS.exists_unitary_leftRegular_fourier}
    \leanok
    \uses{thm:peps_regular_fourier_multiplicities}
    There are positive dimensions $d_i$, pairwise inequivalent irreducible
    unitary matrix representations $D_i$, and an orthonormal basis
    $(e_{i,j,a})$ of $\mathbb C^G$, in row then multiplicity order, in
    which the actual regular permutation matrices are
    \begin{align}
        U_{\mathrm{reg}}(g)\cong
        \bigoplus_i D_i(g)\otimes I_{d_i}.
        \label{eq:peps_regular_unitary_fourier}
    \end{align}
    The matrices $D_i$ and the orthonormal change of basis are derived
    from the finite group. This is the regular block form used in
    \cite[Section~7]{Schuch2010PEPS}.
\end{theorem}
\begin{proof}\leanok
    Take the matrix representation of one irreducible row in its inherited
    orthonormal basis. This gives a unitary irreducible homomorphism $D_i$.
    The common action on every multiplicity row gives the same matrix
    $D_i(g)$ on each copy. The preceding multiplicity identity supplies
    exactly $d_i$ copies. Interchanging the row and copy indices gives
    (\ref{eq:peps_regular_unitary_fourier}).
\end{proof}
